\label{sec:groupoid}
\subsection{Strata, leaves and transverse measures}
Stage 11 set up the wall foliation of input space at depth $l$. Its leaves are the open faces of the activation tiling: cells
(codimension $0$), walls ($1$), and crossings ($\ge2$). Its groupoid algebra has a finite composition series over the strata, with
subquotients $C^*(\text{isotropy})\otimes\mathcal K$:
\begin{itemize}[leftmargin=*]
\item $C_0(\R^j)$ at normal crossings, where the isotropy is abelian;
\item $C^*(\Aff_+(1))$ at bent crossings, where a deep wall meets a wall that feeds it.
\end{itemize}
The Gaussian $\gamma$ induces a transverse measure on every stratum by disintegration. On a stratum $S$ of codimension $j$ it is the
Gaussian surface density on $S$, with the normal space $\R^j$ as local transversal.
\begin{proposition}[the mean is leaf-space integration; stage 11, restated]\label{prop:leafint}
\[
\E f=(2\pi)^{-1/2}\sum_{\mathrm{codim}\,F=1}\kappa_F\,\alpha_F ,
\]
the codimension-one transverse measure paired with the kink cocycle $\kappa_F$.
The cumulant corrections of order $k$ pair the codimension-$j$ transverse measures, $j\le k-1$, with higher kink cocycles. Only bent
chains carry mixed terms (stage 12: Kikuchi level $=$ codimension). The mean, and every correction to the closure, is an integral over
the leaf space.
\end{proposition}

\subsection{The tile groupoid and the network's gap labels}
Let $R_l$ be the relation ``same cell at depth $l$''. It is proper, and $C^*(R_l)\cong\bigoplus_C\mathcal K(L^2(C))$, so $K_0(C^*(R_l))\cong\Z^{(\text{cells})}$.
\begin{proposition}[proved, with the trace normalized by $\gamma$]\label{prop:gap}
The trace induced by the Gaussian transverse measure gives each cell $C$ the weight $\gamma(C)$. Its range on $K_0$ is the
\emph{gap-labelling group}
\[
\mathrm{Gap}_l=\tau_*K_0(C^*(R_l))=\sum_C\Z\,\gamma(C)\ \subset\R .
\]
For every self-adjoint element $H$ of $C^*(R_l)$ (a cell-covariant observable, for instance any function of the gate code), the
integrated density of states $\tau(\1_{(-\infty,E]}(H))$ at a spectral gap lies in $\mathrm{Gap}_l$. Refinement with depth gives
$\mathrm{Gap}_l\subset\mathrm{Gap}_{l+1}$.
\end{proposition}
This is the commutative half of the leaf space. Its K-theoretic ``lattice'' is the module of cell masses, the global code of gate
patterns weighted by the Gaussian. Bellissard's gap labelling for quasicrystals has the same shape, with patch frequencies in place
of cell masses.

\subsection{The non-commutative strata and their modular dilation}
\begin{theorem}[structural; proved at the Lie-algebra level]\label{thm:affiso}
Stage 11 found the isotropy algebra at a bent crossing, $\aff(1)=\langle E,\vartheta\rangle$ with $[E,\vartheta]=\vartheta$, where $E$ is the Euler (dilation) field
and $\vartheta$ the bending shear. It is isomorphic to the Lie algebra of the location--scale group $\Aff_+(1)=\{x\mapsto\sigma x+\mu\}$, which acts
simply transitively on univariate Gaussians. That group is stage 12's hyperbolic upper tier. The isomorphism sends $E$ to the
scale generator $D$ and $\vartheta$ to the translation $P$ ($[D,P]=P$).
\end{theorem}
\begin{theorem}[the modular flow is the dilation; known ingredients, applied]\label{thm:modular}
$\Aff_+(1)$ is not unimodular. Its left Haar measure is $d\sigma\,d\mu/\sigma^2$, its right Haar measure $d\sigma\,d\mu/\sigma$, and its modular
function $\Delta(\sigma,\mu)=\sigma^{-1}$. The Plancherel weight $\varphi$ on the group von Neumann algebra $L(\Aff_+(1))$ is faithful, normal and
semifinite, but not a trace. Its modular group is (Takesaki, Haagerup)
\[
\sigma^\varphi_s\big(\lambda_{(\sigma,\mu)}\big)=\Delta(\sigma,\mu)^{is}\,\lambda_{(\sigma,\mu)}=\sigma^{-is}\,\lambda_{(\sigma,\mu)} .
\]
The modular Hamiltonian $\log\Delta=-\log\sigma$ is the logarithm of the dilation.
\end{theorem}
\begin{corollary}[the critical gain is modular]
On the non-commutative strata, the natural transverse weight is not a trace, and its modular flow is the dilation. The mode that
stage 13 found critical (transport factor $1$, a conserved charge) is the modular Hamiltonian of the network's non-commutative
strata. Three statements are one: the scale is conserved (homogeneity), never frozen by contacts (Proposition~\ref{prop:dilfree}), and
generates the modular flow where the foliation is non-commutative.
\end{corollary}
\begin{remark}[what follows, and what does not]
\begin{itemize}[leftmargin=*]
\item By Connes' Thom isomorphism, the group $C^*$-algebra of $\Aff_+(1)$ has $K_0=\Z$ and $K_1=0$. Gap labels on the bent strata cannot
come from a trace. They need the dilation-twisted (KMS) weight, that is, the twisted cyclic cohomology of Connes and Moscovici.
\item Tomita--Takesaki gives $JL(G)J=R(G)$, so the commutant of the left group algebra is the right one. Stage 11 had forward activations
and backward sensitivities as left- and right-harmonic states. Reading them as the left and right sides of the bent-stratum algebra is
an analogy we have not made precise.
\end{itemize}
\end{remark}

\subsection{An analogy with the space of lattices}
Klartag's process lives, in its dual form, on the space of unimodular lattices $X_n=SL_n(\Z)\backslash SL_n(\R)$, in its Mahler-compact thick
part. The Bost--Connes and Connes--Marcolli systems make a non-commutative space of lattices with the \emph{scaling} as time evolution.
Its KMS states encode the arithmetic, with partition function $\zeta(\beta)$. For the network, the same dilation is the modular flow of the
non-commutative strata. This suggests one well-defined object, the \emph{radial partition function}
\[
Z_l(\beta)=\frac{\E|h_l(x)|^\beta}{\E|x|^\beta}=\int_{S^{n-1}}|h_l(\hat x)|^\beta\,d\sigma(\hat x),
\]
the Mellin transform of the gain. Its logarithm generates the gain's cumulants: $\frac{d^2}{d\beta^2}\log Z|_{0}=\Var\log|h_l(\hat x)|$, the dilation
charge, about $V_l/4$ minus the input radius's share. It is the exact object a dilation-consistent closure must transport (stage 13).
The arithmetic content of Bost--Connes has no counterpart here, and we do not claim one.
